\section{Заключение}

После выполнения лабораторной работы у меня появились ответы на следующие вопросы.

\paragraph{Почему метод Эйлера везде совпал с аналитическим решением?} Подозреваю, что это связано с адаптивным шагом \textit{(зависимости от \( \varepsilon \))} и, как следствие, большим числом узлов. В частности, некоторые примеры изображены с предварительным вычислением значений на миллионе узлов \textit{(хотя мы отрисовываем лишь 100 из них для интерполяционного многочлена)}. 

\paragraph{Какая асимптотическая сложность у алгоритма?} 
Расчёт метода Эйлера на фиксированной сетке из $N$ шагов имеет линейную сложность $O(N)$, так как алгоритм один раз последовательно проходит по всем узлам интервала.

В адаптивной реализации для проверки точности по правилу Рунге сетка каждый раз увеличивается ровно в два раза:
\[
2 + 4 + 8 + \dots + N_{max} = 2N_{max} - 2 = O(N_{max})
\]
где $N_{max}$ — финальное число узлов сетки, на котором разность двух приближений стала меньше $\epsilon$.

То есть итоговая асимптотическая сложность алгоритма адаптивного метода Эйлера линейна и составляет $O(N_{max})$.

\paragraph{В чём преимущество метода Эйлера?} 

Алгоритм невероятно прост: что его блок-схема проста, что программная реализация \textit{(визуализация не в счёт)}. Также расчёт каждого шага занимает крайне малое время --- требуется всего одно вычисление правой части ОДУ. Кроме того, явный метод Эйлера не требует хранения предыстории вычислений в памяти.

\paragraph{Какие недостатки обнаружены у метода Эйлера?} 
Метод имеет первый порядок точности $O(h)$. Из-за этого для достижения высокой точности приходится брать очень-очень малый шаг, что кратно увеличивает число итераций. По идее ещё должна быстро копиться ошибка, но на десяти примерах этого я не заметил. Но однозначно могу отметить, что при отдалении интересуемой точки время расчёта существенно увеличивается.

\paragraph{Когда метод не применим}
В теории метод Эйлера неприменим для решения жёстких систем дифференциальных уравнений, где ограничение на шаг оказывается слишком жёстким для сохранения устойчивости. Однако у меня не возникло проблем для такой системы \textit{(см. №5)}.

\paragraph{Является ли метод устойчивым?} Для начала нужно понять, что такое устойчивый метод.

\textit{Устойчивость численного метода} --- свойство ограничивать рост погрешностей в процессе итерационного вычисления.

Если метод неустойчив, то случайная ошибка на начальных итерациях будет постепенно увеличиваться, что приведет к расходимости численного решения с аналитическим. 

Метод Эйлера является \textit{условно устойчивым}, то есть она зависит от величины шага $h$. И если выбрать $h$ больше некоторого критического, то накопление ошибки приведет к отклонению графика решения от аналитического.